% ============================================================
% 1. СПЕЦИАЛЬНАЯ ЧАСТЬ — 1.1. Объект исследования
% ============================================================
\section{СПЕЦИАЛЬНАЯ ЧАСТЬ}

\subsection{Объект исследования}

\subsubsection{Краткий исторический очерк развития дисциплины}

Сопротивление материалов как научная дисциплина имеет многовековую историю, уходящую корнями в эпоху Возрождения. Первые систематические наблюдения над поведением балок под нагрузкой связаны с именем Галилео Галилея (1564--1642), который в трактате <<Беседы и математические доказательства, касающиеся двух новых наук>> (1638) сформулировал задачу о прочности консольной балки~\cite{timoshenko}. Галилей первым поставил вопрос о том, почему балка одной длины способна выдерживать определённую нагрузку, и предложил количественное описание этой зависимости, хотя его модель распределения напряжений по сечению (равномерное напряжение) впоследствии была признана ошибочной.

Принципиальный прорыв в теории изгиба балок совершили Якоб~Бернулли (1655--1705), впервые связавший кривизну изогнутой оси с изгибающим моментом, и Леонард Эйлер (1707--1783), который в~1744~году получил классическое уравнение изогнутой оси балки, носящее ныне название уравнения Эйлера--Бернулли~\cite{timoshenko, timoshenko1983}:
\[
  EI \, \frac{d^2 v}{dx^2} = M(x).
\]
Это уравнение связало чисто геометрическую характеристику (кривизну деформированной оси) с силовой характеристикой (изгибающим моментом) через изгибную жёсткость~$EI$, открыв путь к количественному инженерному расчёту балочных конструкций.

Вторая половина XIX века была ознаменована работами французского инженера Адемара Барре де Сен-Венана (1797--1886), установившего принцип локальности (известный как принцип Сен-Венана) и давшего строгую постановку задач теории упругости. В этот период были сформированы основные соотношения между внешними нагрузками, внутренними силовыми факторами и деформациями, которые легли в основу современных расчётов балок и стержневых систем~\cite{feodosev, darkov}.

В~XX~веке дисциплина получила дальнейшее развитие благодаря работам С.\,П.~Тимошенко, Н.\,М.~Беляева, В.\,И.~Феодосьева, А.\,Ф.~Смирнова, И.\,А.~Биргера и других выдающихся учёных, заложивших методическую основу преподавания сопротивления материалов в российских и советских технических вузах~\cite{belyaev, feodosev, pisarenko}. С появлением цифровых вычислительных машин в~1950--60-х годах развитие получили методы конечных элементов (МКЭ) и численные подходы, которые, однако, не заменили классические аналитические методы в учебных курсах, а дополнили их.

\subsubsection{Сопротивление материалов как учебная дисциплина}

Сопротивление материалов (сопромат) --- это раздел механики деформируемого твёрдого тела, изучающий напряжения, деформации и условия прочности элементов конструкций~\cite{feodosev, aleksandrov}. Дисциплина занимает центральное место в подготовке инженеров-механиков, строителей, конструкторов транспортных средств и других технических специалистов~\cite{darkov, stepin}. В рамках учебного процесса по направлению 01.03.02 <<Прикладная математика и информатика>> курс сопротивления материалов служит связующим звеном между общими математическими дисциплинами (математический анализ, дифференциальные уравнения) и прикладными задачами компьютерного моделирования поведения конструкций.

Ключевые понятия, осваиваемые студентами в курсе сопротивления материалов~\cite{belyaev, pisarenko}:

\begin{itemize}
  \item \textbf{Напряжения} --- мера внутренних сил, возникающих в теле под действием внешних нагрузок. Нормальное напряжение~$\sigma$ характеризует сопротивление растяжению или сжатию, касательное напряжение~$\tau$ --- сопротивление сдвигу.
  \item \textbf{Деформации} --- изменение формы и размеров тела. Относительная линейная деформация~$\varepsilon$ и относительный сдвиг~$\gamma$ связаны с напряжениями через закон Гука: $\sigma = E \cdot \varepsilon$, где~$E$ --- модуль упругости (модуль Юнга)~\cite{timoshenko1983}.
  \item \textbf{Эпюры внутренних силовых факторов} --- графики распределения изгибающего момента~$M(x)$, поперечной силы~$Q(x)$ и продольной силы~$N(x)$ вдоль оси стержня. Построение эпюр является основной расчётной задачей курса~\cite{gere2018, hibbeler2016}.
  \item \textbf{Условия прочности} --- неравенства, определяющие допустимый уровень напряжений: $\sigma_{\max} \leq [\sigma]$, где~$[\sigma]$ --- допускаемое напряжение материала~\cite{beer2020}.
  \item \textbf{Прогиб и угол поворота} --- перемещения точек оси балки $v(x)$ и наклон касательной $\theta(x)$, определяемые двойным интегрированием уравнения изогнутой оси~\cite{popov1998}.
\end{itemize}

\subsubsection{Анализ трудностей восприятия}

Многочисленные исследования в области инженерного образования фиксируют характерные трудности, с которыми сталкиваются студенты при изучении сопротивления материалов~\cite{philpot2006, philpot2003}:

\begin{enumerate}
  \item \textbf{Абстрактность эпюр.} Студенты часто воспринимают эпюры как формальные графики, не связывая их с физическим процессом деформирования. Переход от расчётной схемы к эпюре требует мысленного <<рассечения>> балки и анализа равновесия отделённой части, что представляет значительную когнитивную нагрузку~\cite{steif2005}.

  \item \textbf{Сложность связи между нагрузкой и деформацией.} Студенты затрудняются объяснить, почему при одинаковой нагрузке балки различных типов деформируются по-разному; почему увеличение длины пролёта на~20\% может увеличить прогиб почти вдвое (зависимость~$v \sim L^3$)~\cite{feodosev}.

  \item \textbf{Ошибки в знаках.} Правила знаков для $Q$ и $M$ --- одна из наиболее частых причин ошибок при построении эпюр. Студенты путают <<растянутое>> и <<сжатое>> волокно, неверно определяют направление реакций~\cite{belyaev, darkov}.

  \item \textbf{Непонимание граничных условий.} Различие между шарнирной опорой, защемлением и свободным концом не всегда осознаётся на физическом уровне, что приводит к ошибкам при определении реакций и постоянных интегрирования~\cite{aleksandrov}.

  \item \textbf{Отсутствие визуальной обратной связи.} При работе с аналитическими формулами студент не видит, как изменение расчётной схемы и нагрузок влияет на форму эпюр и деформированную ось в реальном времени~\cite{balandina, vasilenko}.
\end{enumerate}

Исследование Freeman и др.~\cite{freeman2014} на выборке из 225 работ показало, что активное обучение (active learning) повышает средний балл на экзаменах на~6\% и снижает вероятность получения неудовлетворительной оценки в~1.5~раза по сравнению с традиционными лекциями. Применение интерактивных визуализаций в инженерных дисциплинах рассматривается как один из наиболее эффективных подходов к активному обучению~\cite{makarenko, solovyev}.

\subsubsection{Обоснование выбора статически определимых систем}

В качестве предмета исследования выбраны статически определимые балочные системы по следующим основаниям~\cite{feodosev, aleksandrov}:

\begin{enumerate}
  \item \textbf{Учебная последовательность.} Статически определимые системы изучаются первыми в курсе сопромата и составляют основу для перехода к статически неопределимым задачам.
  \item \textbf{Аналитическая разрешимость.} Реакции опор определяются из уравнений статики без привлечения уравнений совместности деформаций, что делает алгоритм расчёта прозрачным для обучения~\cite{darkov}.
  \item \textbf{Охват ключевых понятий.} На примере трёх типов балок (консольная, шарнирно-опёртая, с вылетом) можно продемонстрировать все основные закономерности: скачки эпюр, экстремумы, связь нагрузки, поперечной силы и изгибающего момента~\cite{belyaev}.
  \item \textbf{Практическая значимость.} Статически определимые системы широко встречаются в практике: мостовые балки, элементы рам, кронштейны, балки перекрытий~\cite{gere2018, beer2020}.
\end{enumerate}

Таким образом, объектом исследования является процесс обучения студентов технических вузов основам сопротивления материалов с использованием интерактивных цифровых средств. Предметом исследования --- интерактивная веб-платформа для визуализации напряжённо-деформированного состояния статически определимых балочных систем.

\subsubsection{Педагогические предпосылки исследования}

С точки зрения педагогического дизайна, разработка обучающей платформы опирается на три теоретических основания.

\textbf{Теория активного обучения} (Active Learning), развитая в работах~\cite{freeman2014, makarenko}, утверждает, что усвоение материала значительно эффективнее при непосредственном взаимодействии обучающегося с изучаемым объектом, а не при пассивном восприятии. В контексте сопротивления материалов это означает, что студент должен не просто читать определения изгибающего момента и наблюдать готовые эпюры, а самостоятельно изменять параметры расчётной схемы и наблюдать следствия своих изменений в реальном времени.

\textbf{Теория когнитивной нагрузки} (Cognitive Load Theory), разработанная Дж.~Свеллером в~1980-х годах, выделяет три типа когнитивной нагрузки: внутреннюю (связанную со сложностью самого материала), постороннюю (связанную с формой представления) и релевантную (направленную на построение ментальных моделей). Грамотно спроектированный интерактивный инструмент снижает постороннюю нагрузку за счёт устранения необходимости вручную выполнять рутинные вычисления и высвобождает когнитивные ресурсы для построения глубокого понимания физических закономерностей.

\textbf{Конструктивистский подход к обучению}, связанный с именем Ж.~Пиаже и Л.\,С.~Выготского, рассматривает обучение как процесс активного построения знания через взаимодействие с предметной средой. В рамках этого подхода важно обеспечить <<зону ближайшего развития>>: задачи должны быть достаточно сложными, чтобы стимулировать мышление, но достаточно доступными, чтобы студент мог их решить при минимальной поддержке. Система контекстных подсказок платформы SopromatLab реализует именно этот принцип, предоставляя объяснения по мере необходимости.

Указанные педагогические основания определяют структуру разрабатываемой платформы и её ключевые отличия от существующих онлайн-калькуляторов, ориентированных исключительно на получение результата расчёта.
